% Generated by docs/generate_user_guide.py. Do not edit directly.
% Source: utilities/descriptions.yaml: metrics; metrics registry

\begin{longtable}{@{}>{\raggedright\arraybackslash}p{0.17\textwidth}>{\raggedright\arraybackslash}p{0.18\textwidth}>{\raggedright\arraybackslash}p{0.20\textwidth}>{\raggedright\arraybackslash}p{0.37\textwidth}@{}}
\toprule
\textbf{Metric} & \textbf{Internal key} & \textbf{Preferred direction} & \textbf{Description} \\
\midrule
\endfirsthead
\toprule
\textbf{Metric} & \textbf{Internal key} & \textbf{Preferred direction} & \textbf{Description} \\
\midrule
\endhead
\(R^2\) & r\_squared & Higher is better & Coefficient of determination. Higher values indicate that the model explains more variance in the observed response. \\
Adjusted \(R^2\) & adjusted\_r\_squared & Higher is better & \(R^2\) adjusted for the number of fitted parameters, which can penalize unnecessary model complexity. \\
MSE & mse & Lower is better & Mean squared error between observed and fitted values. Lower values indicate smaller average squared residuals. \\
RMSE & rmse & Lower is better & Root mean squared error, in the same response units as the fitted data. Lower values indicate smaller typical residuals. \\
MAE & mae & Lower is better & Mean absolute error, in the same response units as the fitted data. Lower values indicate smaller average absolute residuals. \\
NRMSE (mean) & mnrmse & Lower is better & RMSE normalized by the response mean. Lower values indicate smaller error relative to the response magnitude. \\
NRMSE (std. dev.) & nrmse & Lower is better & RMSE normalized by the response standard deviation. Lower values indicate smaller error relative to response variability. \\
AIC & aic & Lower is better & Akaike Information Criterion. Information-criterion metrics balance how well a model fits the data with model complexity (the number of fitted parameters). Among comparable fitted models, a lower AIC value generally indicates a better model representation: the model explains the data well without unnecessary complexity. Absolute AIC values, including negative values, should not be interpreted as standalone quality grades. \\
AIC\textsubscript{c} & aicc & Lower is better & Corrected Akaike Information Criterion. AIC\textsubscript{c} applies a small-sample correction to AIC and is often preferred when the number of fitted observations (n) is limited relative to the number of fitted parameters (k), with n/k < 40 commonly used as a rule of thumb. Among comparable fitted models, a lower AIC\textsubscript{c} value generally indicates a better balance of fit quality and model complexity; AIC\textsubscript{c} is undefined when there are too few observations relative to the number of parameters. \\
BIC & bic & Lower is better & Bayesian Information Criterion. BIC balances how well a model fits the data with model complexity, usually applying a stronger penalty for additional parameters than AIC. Among comparable fitted models, a lower BIC value generally indicates a better model representation without unnecessary complexity. Absolute BIC values, including negative values, should not be interpreted as standalone quality grades. \\
\bottomrule
\end{longtable}
